\begin{tablalafrox}{Diccionario del Producto 2: Requisitos y diseño aprobados. Fuente: según cita por cuenta en cada encabezado.}{tab:T14-dic-prod-2}
  {L{1.2cm} Y Y Y C{0.9cm} L{2.3cm}}
  {\cab{Código} & \cab{Paquete de trabajo} & \cab{Entregable} & \cab{Criterio de aceptación} & \cab{Resp.} & \cab{Período}}
  \multicolumn{6}{>{\columncolor{lafroxBg}\RaggedRight\arraybackslash}p{\dimexpr\textwidth-2\tabcolsep\relax}}{\textbf{Cuenta 2.1 — Requisitos y líneas base de alcance.} Esta cuenta reúne los requisitos y las líneas base de alcance, incluida la captura del conocimiento del planificador y el levantamiento de interfaces. La línea base de la Etapa 1 se aprueba en el H1; la de la Etapa 2 incorpora canal moderno y costo de servir y se aprueba en el H8. {\color{lafroxGris}Fuente: (Bases Administrativas, arts. 17.2–17.6 y Formulario T-14, p. 63); (Bases Técnicas Transversales, secciones 17.2–17.6, pp. 32–34).}} \\
  2.1.1 & Documento de especificación de requerimientos y línea base de alcance de la Etapa 1 (H1) & Especificación de requerimientos aprobada (línea base de alcance). & El 100 \% de los requerimientos de la Etapa 1 tiene módulo y criterio de aceptación, y la línea base está firmada en el mes 2. & ARQ & E1 · Elaboración preliminar: mes 1. Aprobación: H1, mes 2. \\
  2.1.2 & Documento de línea base de alcance de la Etapa 2: canal moderno y costo de servir (H8) & Línea base de alcance de la Etapa 2. & Está aprobada por la Contraparte Técnica en el mes 14. & ARQ & E2 · Elaboración: mes 13. Aprobación: H8, mes 14. \\
  2.1.3 & Matriz de trazabilidad: requerimiento, módulo, paquete y prueba (H1) & Matriz de trazabilidad. & Ningún requerimiento ofertado queda sin paquete ni sin prueba. & CAL & E1 · Mes 2 (H1). Se actualiza con cada cambio aprobado. \\
  2.1.4 & Documento de reglas de ruteo capturadas al planificador de rutas & Catálogo de reglas de ruteo validado por don Hugo. & Don Hugo firma que las reglas representan su criterio. Las rutas que el software genera con esas reglas resultan operables cuando se comparan con las suyas. & IMP & E1 · Meses 1 a 3. \\
  2.1.5 & Documento de especificación de las interfaces sin documentación del ERP, del WMS de 2013 y de la telemetría del tercero, y perfil de intercambio con las cadenas & Especificación inicial de las interfaces (formato, volumen y horario), con el perfil de intercambio por cadena. & Para H2, cada interfaz cuenta con contrato y una prueba de conexión exitosa en el ambiente de Desarrollo. Durante el levantamiento se acuerdan y registran con cada cadena el perfil de mensajes, el acceso de prueba y el calendario de certificación; los de la principal cadena se confirman antes del 1 de octubre de 2027. & ARQ & Ambas · Levantamiento y entrega inicial: 1 al 9 de febrero de 2027. Validación de contratos de interfaz y pruebas de conexión: hasta H2 (mes 4). \\
  \midrule
  \multicolumn{6}{>{\columncolor{lafroxBg}\RaggedRight\arraybackslash}p{\dimexpr\textwidth-2\tabcolsep\relax}}{\textbf{Cuenta 2.2 — Arquitectura de la solución.} Esta cuenta reúne los planos técnicos del software y de su emplazamiento. La Etapa 1 comprende las arquitecturas lógica y física y el modelo de datos; la Etapa 2 añade el diseño detallado de M11 y del intercambio con las cadenas. {\color{lafroxGris}Fuente: (Bases Administrativas, arts. 17.2–17.6 y Formulario T-14, p. 63); (Bases Técnicas Transversales, secciones 17.2–17.6, pp. 32–34).}} \\
  2.2.1 & Documento de arquitectura lógica: capas, módulos M1 a M12 e interfaces & Documento de arquitectura lógica. & Está aprobado en el H2, y cada requerimiento de la línea base apunta a un módulo. & ARQ & E1 · Mes 2 (H2). \\
  2.2.2 & Documento de arquitectura física híbrida: nube, centros de distribución y plataformas de cross-docking & Documento de arquitectura física, con el estudio de cobertura verificado en sitio. & Está aprobado en el H2, y cada componente tiene su lugar y su razón. La verificación en sitio fija los puntos de acceso antes del cableado (SD4, Anexo 4-W.13). & ARQ & E1 · Mes 2 (H2). \\
  2.2.3 & Modelo de datos y diccionario de datos, con la trazabilidad por lote & Modelo de datos y diccionario de datos. & Está aprobado en el H2, y un retiro simulado recorre un lote desde la recepción hasta sus clientes. & DAT & E1 · Mes 2 (H2). \\
  2.2.4 & Documento de diseño detallado de M11 y del intercambio con las cadenas & Documento de diseño detallado de la Etapa 2. & Está aprobado en el H8 y no cambia ningún contrato de la Etapa 1. & ARQ & E2 · Mes 13 (H8). \\
  \midrule
  \multicolumn{6}{>{\columncolor{lafroxBg}\RaggedRight\arraybackslash}p{\dimexpr\textwidth-2\tabcolsep\relax}}{\textbf{Cuenta 2.3 — Diseño de la seguridad.} Esta cuenta de control define cómo se protegen los datos y los accesos. {\color{lafroxGris}Fuente: (Bases Administrativas, arts. 17.2–17.6 y Formulario T-14, p. 63); (Bases Técnicas Transversales, secciones 17.2–17.6, pp. 32–34).}} \\
  2.3.1 & Plan de seguridad de la información & Plan de seguridad. & Está aprobado en el H2, y cada requisito de seguridad del T-12 tiene su control. & SEG & Ambas · Mes 2 (H2). \\
  2.3.2 & Modelado de amenazas por componente y por integración externa, actualizado ante cada cambio & Modelo de amenazas con sus controles. & Cada amenaza alta tiene un control probado o un riesgo aceptado por escrito por el CLIENTE. & SEG & Ambas · Meses 3 a 10, y ante cada cambio de arquitectura. \\
  2.3.3 & Plan de respuesta a incidentes de seguridad, con aviso al CLIENTE en dos horas & Plan de respuesta a incidentes de seguridad. & El plan define la clasificación, el escalamiento, los plazos y los responsables, y se ensaya una vez antes del H5. & SEG & Ambas · Mes 3. \\
  2.3.4 & Matriz de clasificación de la información y matriz de controles & Las dos matrices. & Cada tipo de dato tiene su nivel, y cada nivel tiene sus controles implementados. & SEG & Ambas · Mes 2 (con el H2). \\
  \midrule
  \multicolumn{6}{>{\columncolor{lafroxBg}\RaggedRight\arraybackslash}p{\dimexpr\textwidth-2\tabcolsep\relax}}{\textbf{Cuenta 2.4 — Diseño de la sala técnica y de los sitios.} Esta cuenta reúne los planos de la infraestructura física de la sala técnica y los sitios. La sala actual de Talca tiene 25 m², un aire acondicionado de pared y una UPS que dura 10 minutos, y no cumple el estándar exigido (Caso, cap. 5 y 16). {\color{lafroxGris}Fuente: (Bases Administrativas, arts. 17.2–17.6 y Formulario T-14, p. 63); (Bases Técnicas Transversales, secciones 17.2–17.6, pp. 32–34).}} \\
  2.4.1 & Plano de la sala técnica de Talca, con su cálculo eléctrico y térmico & Plano y memoria de cálculo. & La potencia y la climatización calculadas cubren la carga de diseño con redundancia N+1. & SRE & E1 · Meses 1 y 2. \\
  2.4.2 & Plano de distribución de los racks R01 de servidores y R02 de comunicaciones & Plano de distribución de los racks. & Los servidores y las comunicaciones quedan en racks separados (RT-06.05), y cada equipo tiene doble alimentación (RT-08.04). & SRE & E1 · Meses 1 y 2. \\
  2.4.3 & Plano de los gabinetes de borde de Concepción y de las plataformas de cross-docking & Planos de los gabinetes. & Cada gabinete sostiene su sitio operando sin enlace durante el tiempo declarado. & SRE & E1 · Mes 2. \\
  \midrule
  \multicolumn{6}{>{\columncolor{lafroxBg}\RaggedRight\arraybackslash}p{\dimexpr\textwidth-2\tabcolsep\relax}}{\textbf{Cuenta 2.5 — Continuidad del negocio.} Esta cuenta de control reúne los planes que dicen cómo sigue operando Puelche si la solución, un sitio o un proveedor fallan. Va más allá de la recuperación técnica (3.1.3): incluye contingencias auxiliares manuales cuando el proceso las admite. El despacho de 96 camiones entre 05:30 y 07:00 requiere continuidad tecnológica probada, sin sustitución manual ni interrupción. {\color{lafroxGris}Fuente: (Bases Administrativas, arts. 17.2–17.6 y Formulario T-14, p. 63); (Bases Técnicas Transversales, secciones 17.2–17.6, pp. 32–34).}} \\
  2.5.1 & Plan de continuidad del negocio según ISO 22301, con análisis de impacto y procedimientos manuales & Plan de continuidad del negocio con su análisis de impacto. & Cada proceso identifica continuidad y activación aprobadas por CLIENTE; sólo admite apoyo manual donde corresponda. El despacho crítico exige cero interrupción y flujo tecnológico/DTE probado, independientemente del RTO general. & SRE & E1 · Mes 3. \\
  2.5.2 & Plan de continuidad TIC según ISO/IEC 27031 y plan de recuperación ante desastres & Plan de continuidad TIC y plan de recuperación ante desastres. & Los dos planes están articulados con el plan de continuidad del negocio, y el RTO y el RPO coinciden con los comprometidos. & SRE & E1 · Mes 3. \\
  \midrule
  \multicolumn{6}{>{\columncolor{lafroxBg}\RaggedRight\arraybackslash}p{\dimexpr\textwidth-2\tabcolsep\relax}}{\textbf{Cuenta 2.6 — Experiencia de usuario.} Esta cuenta reúne el diseño de la solución con sus usuarios en preparación de la construcción (BA, art. 26, p. 18). Considera a los preparadores con guantes a −22 °C, preventistas sin señal, conductores externos que se incorporan en minutos y almaceneros sin internet. {\color{lafroxGris}Fuente: (Bases Administrativas, arts. 17.2–17.6 y Formulario T-14, p. 63); (Bases Técnicas Transversales, secciones 17.2–17.6, pp. 32–34).}} \\
  2.6.1 & Investigación con preparadores, preventistas, conductores y almaceneros & Informe de investigación de usuarios por perfil. & Cada perfil operacional tiene sus tareas críticas y sus condiciones de uso documentadas. & IMP & E1 · Mes 2. \\
  2.6.2 & Prototipos y pruebas de usabilidad con participantes del CLIENTE & Prototipos e informe de las pruebas de usabilidad. & Cada flujo crítico se probó con usuarios de su perfil, y los hallazgos están incorporados al diseño antes de construir el módulo (D-05). & IMP & E1 · Mes 3. \\
  2.6.3 & Indicadores de usabilidad comprometidos por transacción crítica & Tabla de indicadores de usabilidad. & Cada transacción crítica tiene sus cuatro indicadores, que se verifican en las pruebas de aceptación (6.2.1). & IMP & E1 · Mes 4. \\
  \midrule
  \multicolumn{6}{>{\columncolor{lafroxBg}\RaggedRight\arraybackslash}p{\dimexpr\textwidth-2\tabcolsep\relax}}{\textbf{Cuenta 2.7 — Sistema de cadena de frío validado.} Esta cuenta reúne la validación del kit de cadena de frío, el mapeo térmico y el piloto en camiones. {\color{lafroxGris}Fuente: (Bases Administrativas, arts. 17.2–17.6 y Formulario T-14, p. 63); (Bases Técnicas Transversales, secciones 17.2–17.6, pp. 32–34).}} \\
  2.7.1 & Validación del diseño del kit en banco de pruebas & Kit completo validado con carga de ensayo & V1–V3, V5–V8, V11, V13, V14, V16, V17 y V19 aprobadas y registradas, sin alimento comercial. & SRE & E1 · Mes 3; antes del mapeo y del piloto. \\
  2.7.2 & Mapeo térmico por tipo de carrocería & Punto crítico, volumen de botellas y tolerancias por tipo de caja & V4 compara núcleo de producto de ensayo, simuladores y aire; Calidad del CLIENTE valida punto y tolerancias con LafroX. & ARQ & E1 · Mes 4; antes del piloto. \\
  2.7.3 & Piloto de cadena de frío en cuatro camiones & Informe de seis semanas de piloto & Dos camiones rurales, uno urbano y uno de transportista; V18, V20 y V21 con evidencia; aprobado antes de la ola de reparto. & IMP & E1 · Meses 10–11; seis semanas antes de la ola 3. \\
  \midrule
  \multicolumn{6}{>{\columncolor{lafroxBg}\RaggedRight\arraybackslash}p{\dimexpr\textwidth-2\tabcolsep\relax}}{\textbf{Cuenta 2.8 — Aprobación de los diseños.} Esta cuenta reúne las actas de aprobación de los diseños. El H2 aprueba la arquitectura, el plan de seguridad y el modelo de datos; el H8 aprueba el diseño detallado de la Etapa 2. {\color{lafroxGris}Fuente: (Bases Administrativas, arts. 17.2–17.6 y Formulario T-14, p. 63); (Bases Técnicas Transversales, secciones 17.2–17.6, pp. 32–34).}} \\
  2.8.1 & Acta de aprobación de la arquitectura, el plan de seguridad y el modelo de datos (H2) & Acta del H2. & Está firmada en el mes 4, sin observaciones abiertas. & ARQ & E1 · Mes 4 (H2). \\
  2.8.2 & Acta de aprobación del diseño detallado de la Etapa 2 (H8) & Acta del H8. & Está firmada en el mes 14. & ARQ & E2 · Mes 13; aprobación en el mes 14 (H8). \\
\end{tablalafrox}
